\documentclass[10pt, a4paper]{article}
\usepackage{preamble}

\title{Algebra II}
\author{Luke Phillips}
\date{August 2025}

\begin{document}

\maketitle

\newpage

\tableofcontents

\newpage

\section{Rings, subrings and fields}

\textit{Informal definition}:
A ring is set of mathematical objects where we can add,
subtract and multiply two elements while remaining in the set.

\subsection{Definition of a ring}
\textbf{Binary operations},
these are operations which take two elements as arguments.

A binary operation on a set $R$ is a mapping of the Cartesian product of two elements of $R$,
$R \times R$ to $R$.
For addition this would be
\[
+(x, y) : R \times R \to R\qquad(\text{with $x, y \in R$}).
\]

A ring $R$ has a binary operation $+$ which makes it into an abelian group.

$R$ is an abelian group with respect to $+$ means that:
\begin{enumerate}[label = \roman*)]
    \item $R$ contains an additive identity element.

    \item Associativity.

    \item For every $x \in R$ there exists a $y \in R$ such that $x + y = y + x = 0$.
    Such that $y$ is unique and generally denoted $-x$.

    \item To make the group abelian we need for $x + y = y + x$ for all $x, y \in R$.
\end{enumerate}

\begin{definition}[Ring]
    A ring $R$ is a set together with two binary operations
    (usually written $+, \cdot$)
    called "addition" and "multiplication",
    such that $+$ makes $R$ into an abelian group,
    and such that the following conditions hold:
    
    \begin{enumerate}[label = \roman*)]
        \item \phantom{}[(multiplicative) Identity]
        There exists an element $1 \in R$ such that $1 \cdot x = x \cdot 1 = x$,
        for all $x \in R$.

        \item \phantom{}[Associativity]
        For any $x, y, z \in R$ we have $x(yz) = (xy)z$.

        \item \phantom{}[Distributivity]
        For any $x, y, z \in R$ we have $x(y + z) = xy + xz$ and $(y + z)x = yx + zx$.
    \end{enumerate}
\end{definition}

If the ring is commutative distributivity is equivalent in both cases.

\subsection{The integers modulo \texorpdfstring{$n$}{}}
A
(binary)
relation $\sim$,
on a set $S$ is said to be an equivalence relation,
if and only if it is reflexive,
symmetric and transitive.
That is,
for all $a, b$ and $c$ in the set $S$ we have:
\begin{enumerate}[label = \roman*)]
    \item 
    $a \sim a$
    (reflexivity).

    \item
    $a \sim b \iff b \sim a$
    (symmetry).

    \item
    If $a \sim b$ and $b \sim c$ then $a \sim c$
    (transitivity).
\end{enumerate}

The equivalence class of $a$ under the given relation $\sim$,
denoted by $[a]$
(or $\bar{a}$)
is defined as
\[
[a] = \{a \in S\,:\, x \sim a\}.
\]
For two $a, b \in S$ with $a \sim b$ we have $[a] = [b]$.
When we write $[a]$ we call $a$ the representative of the class $[a]$.

Given a set $S$ with an equivalence relation $\sim$ as above we may construct another set
\[
S / \sim \coloneqq \{[a]\, :\, a \in S\}
\]
That is,
the set of equivalence classes on the set $S$ with respect to the given equivalence relation $\sim$.


\textbf{The integers modulo} $n$.

Let $n \in \N$ be a natural number.
Consider the set of integer numbers $\Z$ and define the following relation:
two integers $a, b$ are equivalent $a \sim b$ if $n \mid (a - b)$.
The most common notation for this is $a \equiv b \mod n$.

For an integer $a$ we denote $\bar{a}$ the equivalence class of $a$,
that is
\[
\bar{a} = \{x \in \Z\,:\, n\text{ divides } x - a\} = \{x \in \Z\,:\, x \equiv a\mod n\}
\]
since when we divide an integer by $n$ we will always get a remainder among the numbers
\[
0, 1, \dotsc, n - 1.
\]
That is,
every integer belongs to one of the following equivalence classes:
\[
\{\bar{0}, \bar{1}, \dotsc, \overline{n - 1}\},
\]
we will denote this set
(of equivalence classes)
by $\Z / n$.
The elements of $\Z / n$ are called residue classes.

\begin{definition}[Subring]
    A subring $S$ of a ring $R$ is a set $S \subseteq R$ such that
    \begin{enumerate}[label = \roman*)]
        \item $0 \in S, 1 \in S$,

        \item $a, b \in S \implies a + b \in S$,
        
        \item $a, b \in S \implies ab \in S$,
        
        \item $a \in S \implies -a \in S$,
    \end{enumerate}
\end{definition}
This can be summed up as "A subring contains $0$ and $1$ and is closed under addition,
subtraction and multiplication".

\begin{definition}[Field]
    A ring $R$ is called a field if
    
    \begin{enumerate}[label = \roman*)]
        \item $R$ is commutative,

        \item $1 \neq 0$ in $R$
        (i.e. $R$ has at least two elements),


        \item For any $a \in R$ such that $a \neq 0$,
        there exists a $b \in R$ such that $ab = 1$
        ($b$ is called the inverse of $a$).
    \end{enumerate}
\end{definition}

\newpage

\section{Integral domains,
units,
polynomial rings}

\begin{definition}
    A ring $R$ which is commutative,
    has at least two elements
    ($0 \neq 1$),
    and does not have any zero divisors apart from $0$,
    is called an integral domain.
\end{definition}
\textit{Note:
A non zero divisor is a number $b$ such that $ab = 0$ where $b \neq 0$.}

\subsection{Examples and non-examples of integral domains}
\begin{proposition}
    If $R$ is an integral domain,
    then the ring of polynomials $R[x] = \{a_0 + a_1x + \dotsi + a_nx ^ n\mid a_i \in R\}$ is also an integral domain.

    \begin{proof}
        We can clearly see that $R[x]$ is commutative and that $0 \neq 1$
        (as this holds in $R$).

        We need to now show that there are no non-zero divisors in $R[x]$ apart from $0$.
        Assume not,
        then for non-zero polynomials $f(x), g(x) \in R[x]$,
        $f(x)g(x) = 0$.
        Then we get
        \begin{align*}
            f(x) &= a_0 + \dotsc + a_nx ^ n,\quad a_n \neq 0, \\
            g(x) &= b_0 + \dotsc + b_nx ^ m,\quad b_m \neq 0
        \end{align*}
        if we expand $f(x)g(x)$ then the term of the highest degree will be $a_nb_mx ^ {n + m}$.
        Since $f(x)g(x) = 0$ $a_nb_mx ^ {n + m} = 0$ which implies $a_nb_m = 0$ and since $a_n, b_m \in R$ which is an integral domain, $a_n = 0$ or $b_m = 0$.
        Contradiction!
    \end{proof}
\end{proposition}

\subsection{The group of units in a ring}
Let $R$ be a ring.
$a \in R$ is called a unit if there exists a $b \in R$ such that $ab = ba = 1$
($b$ is the inverse of $a$ and $b = a ^ {-1}$).

The inverse is unique as,
if $b_1, b_2 \in R$ such that $ab_1 = b_1a = 1 = ab_2 = b_2a$ then
\[
b_1 = b_11 = b_1(ab_2) = (b_1a)b_2 = 1b_2 = b_2.
\]

The set of all units of $R$ is denoted $R ^ {\times}$
(or $R ^ *$).

$R ^ {\times}$ is always a group,
under the multiplication from the ring $R$:
\[
a, b \in R ^ {\times} \implies ab \in R ^ {\times}
\]
if $c, d \in R$ are elements such that $ac = 1$ and $bd = 1$,
then $(ab)(dc) = abcd = ad = 1 = (dc)(ab)$,
so $ab$ is a unit.
We have an identity element $1$ in the group,
multiplication is associative
(as it is in the ring $R$)
and every element in $R ^ {\times}$ has an inverse by definition.
Thus $R ^ {\times}$ is a group.


\begin{proposition}\label{prep:algII:prop:eleunitiffgcd}
    An element $\bar{a} \in \Z / n$  is a unit if and only if $\gcd(a, n) = 1$.

    \begin{proof}
        Let $d = \gcd(a, n)$.
        Then $d \mid a$ and $d \mid n$.

        Suppose $\bar{a}$ is a unit and let $\bar{a} = \bar{b} ^ {-1}$,
        then $\bar{a}\bar{b} = \bar{1}$ which is equivalent to $ab \equiv 1\mod{n}$.
        Thus $ab = 1 + kn$,
        for some $k \in \Z$.
        Since $d \mid (ab)$ and $d \mid kn$ we must have $d \mid 1$.
        Hence $d = 1$.
        
        Conversely,
        assume that $d = 1$.
        The Euclidean algorithm shows that there exist $x, y \in \Z$ such that 
        \[
        xa + ny = d = 1.
        \]
        Therefore $xa \equiv \mod{n}$,
        which means that $\bar{x}\bar{a} = \bar{1}$ in $\Z / n$.
        Hence $\bar{a}$ is a unit.
    \end{proof}
\end{proposition}

\begin{corollary}\label{prep:algII:cor:zpfieldifprime}
    $\Z / p$ is a field if and only if $p$ is a prime.

    \begin{proof}
        By \autoref{prep:algII:prop:eleunitiffgcd},
        if $p$ is a prime,
        then every element $\bar{1}, \dotsc, \overline{p - 1}$ is a unit,
        so $\Z / p$ is a field.

        Going forwards,
        if $\Z / p$ is a field,
        every non-zero $\bar{a} \in \Z / p$ is a unit,
        so by \autoref{prep:algII:prop:eleunitiffgcd},
        $\gcd(a, o) = 1$ for all $a$ such that $1 \leq a \leq p - 1$.
        This means that $p$ has no divisors between $1$ and $p - 1$ other than $1$,
        so $p$ is prime.
    \end{proof}
\end{corollary}

\begin{proposition}
    The ring $\Z / n$ is an integral domain if and only if $n$ is a prime.

    \begin{proof}
        If $n$ is a prime,
        which \autoref{prep:algII:cor:zpfieldifprime} shows that this means that $\Z / n$ is a field and any field is an integral domain.

        Assume $n$ is not prime,
        then $n = ab$,
        $a, b \in \Z$,
        where $2 \leq a, b \leq n - 1$.
        Then
        \[
        \bar{a}\bar{b} = \bar{n} = \bar{0} \in \Z / n,
        \]
        so $\bar{a}$ and $\bar{b}$ are zero divisors in $\Z / n$,
        and so the ring is not an integral domain.
    \end{proof}
\end{proposition}

\subsection{Polynomials over a field and the division algorithm}
From now on,
$F$ will be a field
(e.g. $\Q, \R$ or $\Z / p$ for $p$ prime).

For $f(x) = a_0 + a_1x + \dotsc + a_nx ^ n \in F[x]$,
we define its degree as
\[
\deg{f} \coloneqq \begin{cases}
    \max\{i\mid a_i \neq 0\} &\text{if } f(x) \neq 0 \\
    -\infty &\text{if } f(x) = 0.
\end{cases}
\]
Then we have the following properties for all $f(x), g(x) \in F[x]$:
\begin{itemize}
    \item $\deg(fg) = \deg{f} + \deg{g}$,
    \item $\deg(f + g) \leq \max\{\deg{f}, \deg{g}\}$.
    Note that if $\deg{f} \neq \deg{g}$ then equality.
\end{itemize}


\textbf{The division algorithm in} $\Z$

\begin{proposition}[Division algorithm]
    Let $f(x), g(x) \in F[x]$ with $g(x) \neq 0$.
    Then there exist unique polynomials $q(x), r(x) \in F[x]$ with $\deg{r} < \deg{g}$ such that
    \[
    f(x) = q(x)g(x) + r(x).
    \]

    \begin{proof}
        (Existence)
        If $\deg{g} > \deg{f}$,
        take $q(x) = 0$ and $r(x) = f(x)$.
        If $\deg{g} \leq \deg{f}$,
        let
        \begin{align*}
            f(x) &= a_0 + \dotsc + a_nx ^ n,\quad a_n \neq 0, \\
            g(x) &= b_0 + \dotsc + b_mx ^ m,\quad b_m \neq 0.
        \end{align*}
        Let $d = m - n \geq 0$.
        Using induction on $d$:
        $m = n$ as $0 = m - n \geq 0$.
        Set $q(x) = a_n / b_m$ and
        \[
        r(x) \coloneqq f(x) - q(x)g(x).
        \]
        Note that $\deg{r} < n = \deg{g}$,
        then we have the required $q(x), r(x)$.

        Assume the existence of $q(x), r(x)$ for all $d < k$,
        for some $k \geq 0$.
        Now consider $d = k$,
        so that $n = m + k$.
        Let
        \[
        f_1(x) \coloneqq f(x) - \frac{a_n}{b_m}x ^ {n - m}g(x).
        \]
        Note $\deg{f_1} < \deg{f}$,
        so by induction,
        there exists polynomials $q_1(x), r(x)$ such that $\deg{r} < \deg{g}$ and
        \[
        f_1(x) = q_1(x)g(x) + r(x).
        \]
        Thus
        \[
        f(x) = f_1(x) + \frac{a_n}{b_m}x ^ {n - m}g(x) = (q_1(x) + \frac{a_n}{b_m}x ^ {n - m})g(x) + r(x)
        \]
        so the result is true for $d = k$,
        and hence for all $d \geq 0$ by induction.

        (Uniqueness)
        Suppose $f(x) = q(x)g(x) + r(x) = Q(x)g(x) + R(x)$,
        with $\deg{r} < \deg{g}$ and $\deg{R} < \deg{g}$.
        This implies $\deg(R - r) < \deg{g}$.
        Rearranging,
        we have
        \begin{equation}\label{prep:algII:eq:eq1}\tag{*}
            R(x) - r(x) = (q(x) - Q(x))g(x),
        \end{equation}
        so by the first basic property,
        we get
        \[
        \deg(R - r) = \deg(q - Q) + \deg{g}.
        \]
        Since $\deg(R - r) < \deg{g}$,
        we must have $\deg(q - Q) < 0$.
        This can only happen if $q(x) = Q(x)$.
        Then by \eqref{prep:algII:eq:eq1} we finally get
        \[
        R(x) = r(x).
        \]
    \end{proof}
\end{proposition}

\newpage

\section{Divisibility and greatest common divisor in a ring}
For this section let $R$ be a commutative ring.

\subsection{The notions of divisibility and of greatest common divisor}
\begin{definition}
    Let $a, b \in R$.
    We say that $a$ divides $b$ in $R$
    ($a \mid b$)
    if there exists an $r \in R$ such that $b = ra$.
\end{definition}

\begin{definition}
    Let $R$ be as above and $a, b \in R$.
    An element $d \in R$ is called a greatest common divisor,
    $\gcd(a, b)$,
    of $a$ and $b$ if
    
    \begin{enumerate}[label = (\roman*)]
        \item $d \mid a$ and $d \mid b$.
        \item If $e \in R$ such that $e \mid a$ and $e \mid b$,
        then $e \mid d$.
    \end{enumerate}
\end{definition}















\end{document}